% !TeX root = main.tex

\chapter*{Apéndice: Rutinas de Wolfram Mathematica para Geometría de la Información}
\addcontentsline{toc}{chapter}{Apéndice: Rutinas de Mathematica}

Este apéndice contiene un archivo completo de Wolfram Mathematica que agrupa todas las rutinas fundamentales de geometría de la información: cálculo de la matriz de información de Fisher, símbolos de Christoffel, tensores de Riemann--Ricci, geodésicas, descomposición en familias exponenciales, divergencias (KL, Bregman), transformada de Legendre y verificación del teorema pitagórico.

\begin{quotation}
	\itshape
	El código está preparado para copiarse y pegarse directamente en un Notebook de Mathematica.  Cada modelo se define como una \emph{Association} con claves \texttt{"function"} y \texttt{"assumptions"}, lo que permite que las funciones se adapten automáticamente a las condiciones de cada problema.
\end{quotation}

\section{Cómo usar este apéndice}
\begin{enumerate}
	\item Copia todo el bloque de código que sigue en un nuevo Notebook de Mathematica.
	\item Evalúa cada celda en orden (Shift+Enter en cada bloque).
	\item Define tu modelo como una Association y pasa las funciones a las utilidades.
\end{enumerate}

\section{Código Mathematica}

\begin{lstlisting}[language=Mathematica, frame=none, numbers=none,
basicstyle=\scriptsize\ttfamily,
keywordstyle=\color{AccentDark}\bfseries,
commentstyle=\color{Accent!60}\itshape,
backgroundcolor=\color{AccentLight!10},
columns=fullflexible,
breaklines=true,
tabsize=2,
xleftmargin=4pt, xrightmargin=4pt,
showstringspaces=false]

(* =================================================== *)
(*  CODIGO UNIFICADO DE GEOMETRIA DE LA INFORMACION    *)
(*  EN MATHEMATICA                                     *)
(* =================================================== *)

(*
Cada modelo se puede definir como una Association con dos claves:
  "function"    -> la funcion de densidad, y
  "assumptions" -> la lista de supuestos.

Si se pasa el modelo como una funcion simple, se utilizan supuestos por defecto.

Se incluyen versiones para:
  (a) Modelos multiparametricos (p.ej., f(x; a, b)).
  (b) Modelos uniparametricos (p.ej., h(x; theta)).

Ademas, se incorporan rutinas para:
  - Matriz de informacion de Fisher.
  - Hessiana y tensor de terceras derivadas.
  - Calculo de los simbolos de Christoffel.
  - Descomposicion en la forma de la familia exponencial.
  - Ecuaciones geodesicas y solucion/graficacion de las mismas.
  - Calculo de tensores (Riemann, Ricci y curvatura escalar).
  - Visualizaciones integrales.
  - Divergencias (KL, Bregman), transformada de Legendre.
  - Verificacion del teorema pitagorico.
*)

(* =================================================== *)
(*   Seccion 1: Rutinas Multiparametricas              *)
(* =================================================== *)

(* 1. Matriz de Informacion de Fisher (Multiparametrica) *)
FisherInformationMatrix[model_, a_, b_, {x_, xmin_, xmax_}] :=
Module[{f, assumptions, L, dLda, dLdb, Iaa, Iab, Ibb},
  If[AssociationQ[model],
    f = model["function"];
    assumptions = model["assumptions"],
    f = model;
    assumptions = {a>0, Element[b,Reals], Element[x,Reals], x>=xmin}
  ];
  L = Log[f[x,a,b]];
  dLda = D[L,a];
  dLdb = D[L,b];
  Iaa = Integrate[f[x,a,b](dLda)^2, {x,xmin,xmax},
                  Assumptions->assumptions];
  Iab = Integrate[f[x,a,b] dLda dLdb, {x,xmin,xmax},
                  Assumptions->assumptions];
  Ibb = Integrate[f[x,a,b](dLdb)^2, {x,xmin,xmax},
                  Assumptions->assumptions];
  Simplify[{{Iaa, Iab},{Iab, Ibb}}]
];

(* 2. Matriz Hessiana *)
HessianMatrix[f_, vars_List] :=
  Table[D[f,vars[[i]],vars[[j]]],
        {i,Length[vars]},{j,Length[vars]}];

(* 3. Tensor de Terceras Derivadas Parciales *)
ThirdOrderTensor[f_, vars_List] :=
  Table[D[f,vars[[i]],vars[[j]],vars[[k]]],
        {i,Length[vars]},{j,Length[vars]},{k,Length[vars]}];

(* 4. Simbolos de Christoffel del Segundo Tipo *)
ChristoffelSymbols[g_, coords_List] :=
Module[{n, gInv, Gamma},
  n = Length[coords];
  gInv = Simplify[Inverse[g]];
  Gamma = Table[
    1/2 Sum[
      gInv[[k,l]] (D[g[[i,l]],coords[[j]]]
                   +D[g[[j,l]],coords[[i]]]
                   -D[g[[i,j]],coords[[l]]]),
      {l,1,n}],
    {i,1,n},{j,1,n},{k,1,n}];
  Simplify[Gamma]
];

(* 5. Descomposicion en Familia Exponencial *)
ExponentialFamilyDecomposition[f_,{x_,a_,b_}] :=
Module[{L,termList,hPart,F,Fterms,Fx,F0,
        candidateTerms,gExpr,pExpr,res},
  L = Expand[Simplify[Log[f[x,a,b]]]];
  termList = If[Head[L]===Plus, List@@L, {L}];
  hPart = Total[Select[termList,
           FreeQ[#,a] && FreeQ[#,b]&]];
  F = Simplify[L - hPart];
  Fterms = If[Head[F]===Plus, List@@F, {F}];
  Fx = Total[Select[Fterms, Not[FreeQ[#,x]]&]];
  F0 = Total[Select[Fterms, FreeQ[#,x]&]];
  candidateTerms = Map[
    Function[t,
      gExpr = Simplify[Coefficient[Coefficient[t,a,0],b,0]];
      If[gExpr===0, gExpr=1];
      pExpr = Simplify[t/gExpr];
      If[FreeQ[gExpr,a] && FreeQ[gExpr,b] && FreeQ[pExpr,x],
        {pExpr, gExpr}, "n/a"]
    ],
    If[Head[Fx]===Plus, List@@Fx, {Fx}]
  ];
  If[MemberQ[candidateTerms,"n/a"],
    "El modelo no es miembro de la familia exponencial",
    res = {
      {"h(x)", Simplify[Exp[hPart]]},
      {"[theta, T(x)]", candidateTerms},
      {"psi(theta)", Simplify[-F0]}
    }; res
  ]
];

(* 6. Ecuaciones Geodesicas *)
GeodesicEquations[g_, coords_List, t_] :=
Module[{n, Gamma, fCoords, eqs},
  n = Length[coords];
  Gamma = ChristoffelSymbols[g, coords];
  fCoords = Map[#[t]&, coords];
  eqs = Table[
    D[fCoords[[i]],{t,2}] +
    Sum[(Gamma[[i,j,k]] /. Thread[coords->fCoords])
         D[fCoords[[j]],t] D[fCoords[[k]],t],
        {j,1,n},{k,1,n}] == 0,
    {i,1,n}];
  Simplify[eqs]
];

(* 7. Solucion y Graficacion de Geodesicas *)
SolveAndPlotGeodesics[g_, coords_List,
    initConds_:Automatic, tmax_:5] :=
Module[{n, t, eqns, sols, plots, defaultInit},
  n = Length[coords];
  defaultInit = {
    {coords[[1]][0]==0, coords[[1]]'[0]==1,
     coords[[2]][0]==0, coords[[2]]'[0]==0},
    {coords[[1]][0]==0, coords[[1]]'[0]==0,
     coords[[2]][0]==0, coords[[2]]'[0]==1}
  };
  If[initConds===Automatic, initConds=defaultInit];
  eqns = GeodesicEquations[g, coords, t];
  sols = Table[
    Module[{sol},
      sol = DSolve[Join[eqns,init],
        Table[coords[[i]][t],{i,1,n}], t];
      If[sol==={} || sol===Null,
        sol = NDSolve[Join[eqns,init],
          Table[coords[[i]],{i,1,n}],{t,0,tmax}];
      ]; sol],
    {init,initConds}];
  plots = Table[
    ParametricPlot[
      Evaluate[Table[coords[[i]][t],{i,1,n}]/.sol[[1]]],
      {t,0,tmax}, PlotRange->All,
      AxesLabel->(ToString/@coords)],
    {sol,sols}];
  {sols, Show[plots,PlotRange->All,
    AxesLabel->(ToString/@coords)]}
];

(* 8. Tensores a partir de la Metrica *)
TensorsFromMetric[g_, coords_List] :=
Module[{n,gamma,Riemann,Ricci,gInv,Rscalar},
  n = Length[coords];
  gamma = ChristoffelSymbols[g, coords];
  Riemann = Table[
    Simplify[
      D[gamma[[a,c,d]],coords[[b]]]
      -D[gamma[[b,c,d]],coords[[a]]]
      +Sum[gamma[[a,c,m]]gamma[[b,m,d]]
           -gamma[[b,c,m]]gamma[[a,m,d]],
           {m,1,n}]],
    {a,1,n},{b,1,n},{c,1,n},{d,1,n}];
  Ricci = Table[
    Simplify[Sum[Riemann[[k,j,i,k]],{k,1,n}]],
    {i,1,n},{j,1,n}];
  gInv = Simplify[Inverse[g]];
  Rscalar = Simplify[Sum[gInv[[i,j]]Ricci[[i,j]],
    {i,1,n},{j,1,n}]];
  {{"Tensor de Riemann",Riemann},
   {"Tensor de Ricci",Ricci},
   {"Curvatura Escalar",Rscalar}}
];

(* =================================================== *)
(*   Seccion 2: Rutinas Uniparametricas                *)
(* =================================================== *)

(* 1. Informacion de Fisher 1D *)
FisherInformation1D[model_, param_, {x_, xmin_, xmax_}] :=
Module[{f,assumptions,L,dLdparam,info},
  If[AssociationQ[model],
    f = model["function"];
    assumptions = model["assumptions"],
    f = model;
    assumptions = {param>0, Element[x,Reals], x>=xmin}
  ];
  L = Log[f[x,param]];
  dLdparam = D[L,param];
  info = Integrate[f[x,param](dLdparam)^2,
    {x,xmin,xmax}, Assumptions->assumptions];
  Simplify[info]
];

(* 2. Hessiana 1D *)
Hessian1D[f_,theta_] := D[f,{theta,2}];

(* 3. Tensor de Terceras Derivadas 1D *)
ThirdOrderTensor1D[f_,theta_] := D[f,{theta,3}];

(* 4. Simbolos de Christoffel 1D *)
ChristoffelSymbols1D[g_,theta_] :=
Module[{gInv,gamma},
  gInv = Simplify[1/g];
  gamma = Simplify[(1/2) gInv D[g,theta]];
  gamma
];

(* 5. Descomposicion en Familia Exponencial 1D *)
ExponentialFamilyDecomposition1D[f_,{x_,theta_}] :=
Module[{L,termList,hPart,F,Fterms,Fx,F0,
        candidateTerm,gExprLocal,pExprLocal,res},
  L = Expand[Simplify[Log[f[x,theta]]]];
  termList = If[Head[L]===Plus, List@@L, {L}];
  hPart = Total[Select[termList, FreeQ[#,theta]&]];
  F = Simplify[L - hPart];
  Fterms = If[Head[F]===Plus, List@@F, {F}];
  Fx = Total[Select[Fterms, Not[FreeQ[#,x]]&]];
  F0 = Total[Select[Fterms, FreeQ[#,x]&]];
  candidateTerm = Module[{gEL,pEL},
    gEL = Simplify[Coefficient[Fx,theta,0]];
    If[gEL===0, gEL=1];
    pEL = Simplify[Fx/gEL];
    If[FreeQ[gEL,theta] && FreeQ[pEL,x],
      {pEL,gEL}, "n/a"]
  ];
  If[candidateTerm==="n/a",
    "El modelo no es miembro de la familia exponencial",
    res = {
      {"h(x)", Simplify[Exp[hPart]]},
      {"[theta, T(x)]", candidateTerm},
      {"psi(theta)", Simplify[-F0]}
    }; res
  ]
];

(* 6. Ecuacion Geodesica 1D *)
GeodesicEquations1D[g_,theta_,t_] :=
Module[{Gamma,thetaFunc,eq},
  Gamma = ChristoffelSymbols1D[g,theta];
  thetaFunc = theta[t];
  eq = Simplify[D[thetaFunc,{t,2}] +
    (Gamma /. theta->thetaFunc)(D[thetaFunc,t])^2 == 0];
  eq
];

(* 7. Solucion y Graficacion de Geodesicas 1D *)
SolveAndPlotGeodesics1D[g_,theta_,
    initConds_:Automatic, tmax_:5] :=
Module[{t,eq,sol,plot,defaultInit},
  eq = GeodesicEquations1D[g,theta,t];
  defaultInit = {theta[0]==1, theta'[0]==0.5};
  If[initConds===Automatic, initConds=defaultInit];
  sol = DSolve[Join[{eq},initConds],theta[t],t];
  If[sol==={} || sol===Null,
    sol = NDSolve[Join[{eq},initConds],
      theta,{t,0,tmax}];
    plot = Plot[Evaluate[theta[t]/.sol[[1]]],{t,0,tmax},
      AxesLabel->{"t",ToString[theta]},
      PlotLabel->"Geodesica 1D"],
    plot = Plot[Evaluate[theta[t]/.sol[[1]]],{t,0,tmax},
      AxesLabel->{"t",ToString[theta]},
      PlotLabel->"Geodesica 1D (analitica)"]
  ];
  {sol,plot}
];

(* 8. Tensores a partir de la Metrica 1D *)
(* En 1D el tensor de Riemann se anula por antisimetria:
   R^1_{111} = 0 siempre. La curvatura escalar es 0. *)
TensorsFromMetric1D[g_,theta_] :=
Module[{gInv,gamma},
  gInv = Simplify[1/g];
  gamma = ChristoffelSymbols1D[g,theta];
  {{"Tensor de Riemann", 0},
   {"Tensor de Ricci", 0},
   {"Curvatura Escalar", 0}}
];

(* =================================================== *)
(*   Seccion 3: Funciones de Utilidad Comunes          *)
(* =================================================== *)

(* 1. Informacion de Fisher para distribuciones comunes *)
FisherBernoulli[p_?NumericQ] := Simplify[1/(p(1-p))];

FisherNormal[mu_, sigma_?NumericQ] :=
  Simplify[{{1/sigma^2,0},{0,2/sigma^2}}];

FisherPoisson[lambda_?NumericQ] := Simplify[1/lambda];

FisherExponential[lambda_?NumericQ] := Simplify[1/lambda^2];

(* 2. Divergencia KL numerica *)
KLDivergence[f_,g_,{x_,xmin_,xmax_}] :=
Module[{integrand,assumptions},
  assumptions = {Element[x,Reals],x>=xmin};
  integrand = f Log[f/g];
  Simplify[Integrate[integrand,{x,xmin,xmax},
                     Assumptions->assumptions]]
];

(* 3. Divergencia de Bregman generica *)
BregmanDivergence[F_,p_,q_] :=
Module[{gradF},
  gradF = D[F,#]&;
  (F /. q->p) - (F /. q->q) - gradF[q](p-q)
];

(* 4. Transformada de Legendre-Fenchel simbolica *)
LegendreTransform[F_,x_,eta_] :=
Module[{sol,starFunc},
  sol = Solve[D[F,x]==eta, x];
  If[sol==={} || sol===Null,
    "No se puede resolver para eta = dF/dx",
    sol = First[sol];
    starFunc = (x eta - F) /. sol;
    Simplify[starFunc]
  ]
];

(* 5. Verificacion numerica del teorema pitagorico *)
VerifyPythagoreanKL[model_,thetaTrue_,
                    thetaProj_,thetaQ_] :=
Module[{d1,d2,d3},
  d1 = KLDivergence[model[thetaTrue],
                     model[thetaProj],
                     {x,-Infinity,Infinity}];
  d2 = KLDivergence[model[thetaProj],
                     model[thetaQ],
                     {x,-Infinity,Infinity}];
  d3 = KLDivergence[model[thetaTrue],
                     model[thetaQ],
                     {x,-Infinity,Infinity}];
  Print["D_KL(p||q*) = ",N[d1]];
  Print["D_KL(q*||q) = ",N[d2]];
  Print["Suma       = ",N[d1+d2]];
  Print["D_KL(p||q)  = ",N[d3]];
  Print["Pitagorico: ",
    N[Abs[d3-(d1+d2)]] < 1*^-10]
];

(* 6. Coordenadas duales para familias exponenciales *)
DualCoordinatesBernoulli[p_] := {Log[p/(1-p)],p};
DualCoordinatesPoisson[lambda_] := {Log[lambda],lambda};
DualCoordinatesNormal[mu_,sigma_] := {mu/sigma^2,mu};

(* =================================================== *)
(*              EJEMPLOS DE USO                        *)
(* =================================================== *)

(* --- Ejemplo Multiparametrico --- *)
expModel2D = <|
  "function" -> (a Exp[-a(x-b)] UnitStep[x-b]),
  "assumptions" -> {a>0, Element[b,Reals],
                    Element[x,Reals], x>=b}
|>;
FisherMat = FisherInformationMatrix[expModel2D,
  a, b, {x, b, Infinity}];
Print["Matriz de Fisher (Multiparametrica):"];
Print[FisherMat];

(* --- Ejemplo Uniparametrico --- *)
expModel1D = <|
  "function" -> (a Exp[-a*x] UnitStep[x]),
  "assumptions" -> {a>0, Element[x,Reals], x>=0}
|>;
Print["Fisher (Uniparametrica):"];
Print[FisherInformation1D[expModel1D, a,
  {x, 0, Infinity}]];

(* --- Verificacion rapida --- *)
(* FisherBernoulli[0.5]   deve devolver 4 *)
(* FisherPoisson[3.0]     deve devolver 1/3 *)
(* DualCoordinatesBernoulli[0.7] = {Logit(0.7), 0.7} *)
\end{lstlisting}

\section{Notas sobre las convenciones utilizadas}

\begin{itemize}
	\item \textbf{Notación de los símbolos de Christoffel:} $\Gamma^{k}_{ij}$ se almacenan como \texttt{Gamma[[i,j,k]]} (dos índices inferiores, uno superior).
	\item \textbf{Tensor de Riemann:} En 2D, la convención utilizada es $R^{d}_{abc} = \partial_b \Gamma^{d}_{ac} - \partial_a \Gamma^{d}_{bc} + \Gamma^{d}_{am}\Gamma^{m}_{bc} - \Gamma^{d}_{bm}\Gamma^{m}_{ac}$.  Esto es consistente con las fórmulas de Ricci y curvatura escalar calculadas internamente.
	\item \textbf{1D:} En una dimensión, el tensor de Riemann se anula por antisimetría ($R^{1}_{111}=0$), por lo que la curvatura escalar es siempre cero.  Esto refleja que toda variedad unidimensional es localmente plana.
	\item \textbf{Modelos con Association:} Las funciones aceptan tanto un modelo dado como una Association \texttt{"function" $\to$ ..., "assumptions" $\to$ ...} para mayor flexibilidad.
\end{itemize}
